% GENERATED FILE. Edit canonical lessons, then run npm run generate:books.
% canonical-source-hash: fnv1a64:f567d6edf004445c
% canonical-lessons: ZH-C13-hear-da, ZH-W13-da, ZH-C13-da, ZH-C13-hear-dajia, ZH-C13-dajia, ZH-C13-hear-xue, ZH-W13-xue, ZH-C13-xue, ZH-C13-hear-xuesheng, ZH-W13-sheng, ZH-C13-xuesheng, ZH-C13-hear-daxue, ZH-C13-daxue, ZH-C13-hear-daxuesheng, ZH-C13-daxuesheng, ZH-R13-school-six, ZH-C13-practice

\chapter{From study to university student}
\label{ch:zh-school-six}
\hlchaptermodality{13}

\begin{chapteropening}
\textbf{By the end of this chapter:} \emph{I can understand, say, read and write six HSK Level 1 size-and-school words, gaining \zh{大}, \zh{学} and \zh{生} one gently staged character at a time before recombining them as \zh{大家}, \zh{学生}, \zh{大学} and \zh{大学生}.}

Three newly secured characters produce six useful words, each checked separately in listening, speaking, unassisted reading and no-model writing.
\end{chapteropening}

\section[dà]{Hear \emph{dà} — big}

\label{lesson:ZH-C13-hear-da}



\begin{quote}
Say \textbf{home} and \textbf{family members}.
\end{quote}

\subsection*{What to know first}
\begin{quote}
\emph{dà} — \textbf{big; large}
\end{quote}

Let the voice fall firmly on the single syllable: \emph{dà}. Keep the character covered for now. The next lesson gives its three strokes one at a time.

\begin{tcolorbox}[breakable,colback=teal!4,colframe=teal!35!black,title={Before you move on}]
Source: \href{https://hsk.cn-bj.ufileos.com/3.0/\%E6\%96\%B0\%E7\%89\%88HSK\%E8\%80\%83\%E8\%AF\%95\%E5\%A4\%A7\%E7\%BA\%B21219.pdf}{Official HSK 3.0 syllabus, Level 1 vocabulary}.
\end{tcolorbox}
\section[dà]{\zh{大} — one bar, two falls}

\label{lesson:ZH-W13-da}



\begin{quote}
Say \textbf{big} in Mandarin.
\end{quote}

\subsection*{Script — the shape}
\begin{quote}
\textbf{\zh{一}} + \textbf{\zh{丿}} + \textbf{\zh{㇏}} $\to$ \textbf{\zh{大}}
\end{quote}

Write the horizontal first. Start the left-falling stroke above it and cross through it. Lift once more, return to the crossing, and fall long to the right.

\subsection*{Your turn}
Trace \textbf{\zh{大}} once. Copy it once while saying \textbf{bar, left, right}.

\begin{tcolorbox}[breakable,colback=teal!4,colframe=teal!35!black,title={Before you move on}]
Source: \href{https://raw.githubusercontent.com/chanind/hanzi-writer-data/68d10a4b21150cae5e1ebbd223eed289cf32d90c/data/\%E5\%A4\%A7.json}{Hanzi Writer Data \zh{大}, pinned stroke paths and medians}.
\end{tcolorbox}
\section[dà]{\zh{大} — big on the page}

\label{lesson:ZH-C13-da}



\begin{quote}
Say \textbf{big}, then point to the first stroke of \zh{大}.
\end{quote}

\begin{cousinweb}
\begin{quote}
\textbf{\zh{大}} — \emph{dà} — \textbf{big; large}
\end{quote}

Read the character without pinyin. Then say it once with a clear fourth-tone fall.
\end{cousinweb}

\subsection*{Your turn}
Copy \textbf{\zh{大}} once from the model. Keep the crossing near the centre.

\begin{tcolorbox}[breakable,colback=teal!4,colframe=teal!35!black,title={Before you move on}]
Source: \href{https://hsk.cn-bj.ufileos.com/3.0/\%E6\%96\%B0\%E7\%89\%88HSK\%E8\%80\%83\%E8\%AF\%95\%E5\%A4\%A7\%E7\%BA\%B21219.pdf}{Official HSK 3.0 syllabus, Level 1 vocabulary}.
\end{tcolorbox}
\section[dàjiā]{Hear \emph{dàjiā} — everyone}

\label{lesson:ZH-C13-hear-dajia}



\begin{quote}
Say \emph{dà}, then \emph{jiā}.
\end{quote}

\subsection*{What to know first}
\begin{quote}
\emph{dàjiā} — \textbf{everyone; everybody}
\end{quote}

Hear two beats: falling \emph{dà}, then level \emph{jiā}. The whole word means everyone. Its written form will contain only characters you already know.

\begin{tcolorbox}[breakable,colback=teal!4,colframe=teal!35!black,title={Before you move on}]
Source: \href{https://hsk.cn-bj.ufileos.com/3.0/\%E6\%96\%B0\%E7\%89\%88HSK\%E8\%80\%83\%E8\%AF\%95\%E5\%A4\%A7\%E7\%BA\%B21219.pdf}{Official HSK 3.0 syllabus, Level 1 vocabulary}.
\end{tcolorbox}
\section[dàjiā]{\zh{大家} — everyone on the page}

\label{lesson:ZH-C13-dajia}



\begin{quote}
Write \textbf{\zh{大}}, then \textbf{\zh{家}} with a small gap.
\end{quote}

\begin{cousinweb}
\begin{quote}
\textbf{\zh{大家}} — \emph{dàjiā} — \textbf{everyone}
\end{quote}

Close the gap. Read the two characters as one word and keep both tone shapes.
\end{cousinweb}

\subsection*{Your turn}
Look at \textbf{\zh{大家}}, cover it, then write it once from memory.

\begin{tcolorbox}[breakable,colback=teal!4,colframe=teal!35!black,title={Before you move on}]
Source: \href{https://hsk.cn-bj.ufileos.com/3.0/\%E6\%96\%B0\%E7\%89\%88HSK\%E8\%80\%83\%E8\%AF\%95\%E5\%A4\%A7\%E7\%BA\%B21219.pdf}{Official HSK 3.0 syllabus, Level 1 vocabulary}.
\end{tcolorbox}
\section[xué]{Hear \emph{xué} — study, learn}

\label{lesson:ZH-C13-hear-xue}



\begin{quote}
Say \textbf{everyone}.
\end{quote}

\subsection*{What to know first}
\begin{quote}
\emph{xué} — \textbf{study; learn}
\end{quote}

Begin low enough to rise clearly through \emph{xué}. Stay with sound and meaning; the eight-stroke character comes next in small groups.

\begin{tcolorbox}[breakable,colback=teal!4,colframe=teal!35!black,title={Before you move on}]
Source: \href{https://hsk.cn-bj.ufileos.com/3.0/\%E6\%96\%B0\%E7\%89\%88HSK\%E8\%80\%83\%E8\%AF\%95\%E5\%A4\%A7\%E7\%BA\%B21219.pdf}{Official HSK 3.0 syllabus, Level 1 vocabulary}.
\end{tcolorbox}
\section[xué]{\zh{学} — top, cover, child}

\label{lesson:ZH-W13-xue}



\begin{quote}
Say \textbf{study}, then write \zh{子}.
\end{quote}

\subsection*{Script — the shape}
\begin{quote}
\textbf{\zh{⺍}} + \textbf{\zh{冖}} + \textbf{\zh{子}} $\to$ \textbf{\zh{学}}
\end{quote}

Place the three short marks across the top. Add the covering dot and hooked line. Finish with the familiar three strokes of \zh{子} below. Eight strokes total; do not let the top marks crowd the child.

\subsection*{Your turn}
Trace \textbf{\zh{学}} once. Copy it once while saying \textbf{top, cover, child}.

\begin{tcolorbox}[breakable,colback=teal!4,colframe=teal!35!black,title={Before you move on}]
Source: \href{https://raw.githubusercontent.com/chanind/hanzi-writer-data/68d10a4b21150cae5e1ebbd223eed289cf32d90c/data/\%E5\%AD\%A6.json}{Hanzi Writer Data \zh{学}, pinned stroke paths and medians}.
\end{tcolorbox}
\section[xué]{\zh{学} — study on the page}

\label{lesson:ZH-C13-xue}



\begin{quote}
Say \textbf{study}, then point to \zh{子} inside \zh{学}.
\end{quote}

\begin{cousinweb}
\begin{quote}
\textbf{\zh{学}} — \emph{xué} — \textbf{study; learn}
\end{quote}

Read it without pinyin, then give the syllable a clear rise.
\end{cousinweb}

\subsection*{Your turn}
Copy \textbf{\zh{学}} once. Pause mentally after the top, after the cover, and after \zh{子}.

\begin{tcolorbox}[breakable,colback=teal!4,colframe=teal!35!black,title={Before you move on}]
Source: \href{https://hsk.cn-bj.ufileos.com/3.0/\%E6\%96\%B0\%E7\%89\%88HSK\%E8\%80\%83\%E8\%AF\%95\%E5\%A4\%A7\%E7\%BA\%B21219.pdf}{Official HSK 3.0 syllabus, Level 1 vocabulary}.
\end{tcolorbox}
\section[xuésheng]{Hear \emph{xuésheng} — student}

\label{lesson:ZH-C13-hear-xuesheng}



\begin{quote}
Say \textbf{study}.
\end{quote}

\subsection*{What to know first}
\begin{quote}
\emph{xuésheng} — \textbf{student}
\end{quote}

Rise on \emph{xué}. Keep \emph{sheng} short and light. Hold the second character back; the next lesson gives its five strokes.

\begin{tcolorbox}[breakable,colback=teal!4,colframe=teal!35!black,title={Before you move on}]
Source: \href{https://hsk.cn-bj.ufileos.com/3.0/\%E6\%96\%B0\%E7\%89\%88HSK\%E8\%80\%83\%E8\%AF\%95\%E5\%A4\%A7\%E7\%BA\%B21219.pdf}{Official HSK 3.0 syllabus, Level 1 vocabulary}.
\end{tcolorbox}
\section[sheng]{\zh{生} — five measured strokes}

\label{lesson:ZH-W13-sheng}



\begin{quote}
Say \textbf{student}.
\end{quote}

\subsection*{Script — the shape}
Begin with the short left-fall. Add the upper horizontal and the shorter middle horizontal. Draw the vertical through their centres, then anchor everything with the long bottom horizontal. Five strokes; four lifts.

\subsection*{Your turn}
Trace \textbf{\zh{生}} once. Copy it once while saying \textbf{fall, two bars, through, base}.

\begin{tcolorbox}[breakable,colback=teal!4,colframe=teal!35!black,title={Before you move on}]
Source: \href{https://raw.githubusercontent.com/chanind/hanzi-writer-data/68d10a4b21150cae5e1ebbd223eed289cf32d90c/data/\%E7\%94\%9F.json}{Hanzi Writer Data \zh{生}, pinned stroke paths and medians}.
\end{tcolorbox}
\section[xuésheng]{\zh{学生} — student on the page}

\label{lesson:ZH-C13-xuesheng}



\begin{quote}
Point to \zh{学}, then \zh{生}. Say \textbf{study}, then \textbf{student}.
\end{quote}

\begin{cousinweb}
\begin{quote}
\textbf{\zh{学生}} — \emph{xuésheng} — \textbf{student}
\end{quote}

Read the two-character word without pinyin. Keep the second syllable light.
\end{cousinweb}

\subsection*{Your turn}
Copy \textbf{\zh{学生}} once. Finish all of \zh{学} before beginning \zh{生}.

\begin{tcolorbox}[breakable,colback=teal!4,colframe=teal!35!black,title={Before you move on}]
Source: \href{https://hsk.cn-bj.ufileos.com/3.0/\%E6\%96\%B0\%E7\%89\%88HSK\%E8\%80\%83\%E8\%AF\%95\%E5\%A4\%A7\%E7\%BA\%B21219.pdf}{Official HSK 3.0 syllabus, Level 1 vocabulary}.
\end{tcolorbox}
\section[dàxué]{Hear \emph{dàxué} — university}

\label{lesson:ZH-C13-hear-daxue}



\begin{quote}
Say \textbf{big}, then \textbf{study}.
\end{quote}

\subsection*{What to know first}
\begin{quote}
\emph{dàxué} — \textbf{university; college}
\end{quote}

Let \emph{dà} fall, then let \emph{xué} rise. Together they name a university or college.

\begin{tcolorbox}[breakable,colback=teal!4,colframe=teal!35!black,title={Before you move on}]
Source: \href{https://hsk.cn-bj.ufileos.com/3.0/\%E6\%96\%B0\%E7\%89\%88HSK\%E8\%80\%83\%E8\%AF\%95\%E5\%A4\%A7\%E7\%BA\%B21219.pdf}{Official HSK 3.0 syllabus, Level 1 vocabulary}.
\end{tcolorbox}
\section[dàxué]{\zh{大学} — university on the page}

\label{lesson:ZH-C13-daxue}



\begin{quote}
Read \zh{大} and \zh{学} separately.
\end{quote}

\begin{cousinweb}
\begin{quote}
\textbf{\zh{大学}} — \emph{dàxué} — \textbf{university; college}
\end{quote}

Now read the pair as one word: fall, then rise.
\end{cousinweb}

\subsection*{Your turn}
Look at \textbf{\zh{大学}}, cover it, and write it once from memory.

\begin{tcolorbox}[breakable,colback=teal!4,colframe=teal!35!black,title={Before you move on}]
Source: \href{https://hsk.cn-bj.ufileos.com/3.0/\%E6\%96\%B0\%E7\%89\%88HSK\%E8\%80\%83\%E8\%AF\%95\%E5\%A4\%A7\%E7\%BA\%B21219.pdf}{Official HSK 3.0 syllabus, Level 1 vocabulary}.
\end{tcolorbox}
\section[dàxuésheng]{Hear \emph{dàxuésheng} — university student}

\label{lesson:ZH-C13-hear-daxuesheng}



\begin{quote}
Say \textbf{university}, then \textbf{student}.
\end{quote}

\subsection*{What to know first}
\begin{quote}
\emph{dàxuésheng} — \textbf{university student}
\end{quote}

Say three beats: falling \emph{dà}, rising \emph{xué}, light \emph{sheng}. Do not repeat \emph{xué}; the two familiar words overlap there.

\begin{tcolorbox}[breakable,colback=teal!4,colframe=teal!35!black,title={Before you move on}]
Source: \href{https://hsk.cn-bj.ufileos.com/3.0/\%E6\%96\%B0\%E7\%89\%88HSK\%E8\%80\%83\%E8\%AF\%95\%E5\%A4\%A7\%E7\%BA\%B21219.pdf}{Official HSK 3.0 syllabus, Level 1 vocabulary}.
\end{tcolorbox}
\section[dàxuésheng]{\zh{大学生} — university student on the page}

\label{lesson:ZH-C13-daxuesheng}



\begin{quote}
Read \textbf{\zh{大学}}, then \textbf{\zh{学生}}.
\end{quote}

\begin{cousinweb}
\begin{quote}
\textbf{\zh{大学生}} — \emph{dàxuésheng} — \textbf{university student}
\end{quote}

Track three characters and three beats. The middle \zh{学} belongs to the whole word only once.
\end{cousinweb}

\subsection*{Your turn}
Look at \textbf{\zh{大学生}}, cover it, and write it once. Check one character at a time.

\begin{tcolorbox}[breakable,colback=teal!4,colframe=teal!35!black,title={Before you move on}]
Source: \href{https://hsk.cn-bj.ufileos.com/3.0/\%E6\%96\%B0\%E7\%89\%88HSK\%E8\%80\%83\%E8\%AF\%95\%E5\%A4\%A7\%E7\%BA\%B21219.pdf}{Official HSK 3.0 syllabus, Level 1 vocabulary}.
\end{tcolorbox}
\section[Practice]{Retrieve six size and school words}

\label{lesson:ZH-R13-school-six}



\begin{quote}
Write \zh{大}, \zh{学}, and \zh{生} once each from memory.
\end{quote}

\subsection*{Your turn}
Hear the six words in shuffled order and point to meanings. Then say each word from a meaning card, read each character card, and write only the missed items. Repair one word at a time.

\begin{tcolorbox}[breakable,colback=teal!4,colframe=teal!35!black,title={Before you move on}]
Source: \href{https://hsk.cn-bj.ufileos.com/3.0/\%E6\%96\%B0\%E7\%89\%88HSK\%E8\%80\%83\%E8\%AF\%95\%E5\%A4\%A7\%E7\%BA\%B21219.pdf}{Official HSK 3.0 syllabus, Level 1 vocabulary}.
\end{tcolorbox}
\section[Practice]{Payoff — six school words in four skills}

\label{lesson:ZH-C13-practice}



\begin{quote}
Write \zh{大}, \zh{学}, and \zh{生} without a model.
\end{quote}

\subsection*{Your turn}
1. \textbf{Listening:} hear six words in mixed order and choose every meaning. 2. \textbf{Speaking:} produce all six from meaning cards, with tones audible. 3. \textbf{Reading:} read all six character cards without pinyin. 4. \textbf{Writing:} hear all six and write them without a model.

Pass each skill separately. Repair and retry only the failed skill. This is a school vocabulary checkpoint, not yet a full conversation about education.

\begin{tcolorbox}[breakable,colback=teal!4,colframe=teal!35!black,title={Before you move on}]
Source: \href{https://hsk.cn-bj.ufileos.com/3.0/\%E6\%96\%B0\%E7\%89\%88HSK\%E8\%80\%83\%E8\%AF\%95\%E5\%A4\%A7\%E7\%BA\%B21219.pdf}{Official HSK 3.0 syllabus, Level 1 task outline and vocabulary}.
\end{tcolorbox}
